\section{The Celestial \& Super-Celestial}

\textbf{Exerpts from ``Disputation Against the Judgement of the Astrologers''}

\begin{quotex}
Now because things always proceed in the same order, it is not by chance that there is some other determined cause outside matter and above the natural principle of agency, by whose intention individual things lead appropriately to the most perfect end. Unless one thing is so, the other cannot be.
\end{quotex}


Marsilio Ficino\footnote{\url{http://en.wikipedia.org/wiki/Marsilio\_Ficino}} attacked the New Agers of his day; however, as we shall see, his aim was twofold -- to establish Christianity against them, but to do so by incorporating truly Platonic insights which existed in potential already \emph{in the very material which the astrologers were mutilating}. In other words, Ficino was adding to the tradition, rather than being a parasite (the condition that so many today embrace, whether ``pro-tradition'' or ``anti-tradition''). There is no doubt this is a fine line, or as Christ would say, a ``narrow gate''.

\begin{quotex}
Now in the rational order of the heavens there is no tendency to defect, whereas under the heavens there exists in rational beings the tendency to defect of the will, but this does not derive from the rational part of their nature. In these beings the defect is brought back by God to the order of the good, that is, either into the order of justice through punishment or into the order of compassion through forgiveness.
\end{quotex}

In other words, there is such a thing as the Law of God, or \emph{Lex Aeterna} (and indeed, always will be). Gurdjieff speaks of attaining an astral body which allows one to transcend or ``jump over the moon''. Bondarev (via Steiner) speaks of Yahweh's countenance being manifested under the lunar sign, that is, as a preparation to forming the astral body. The paideia\footnote{\url{http://en.wikipedia.org/wiki/Paideia}} of God has to occur prior as a formation or pre-evangelium; men have to have souls formed before their souls can be saved. Ficino's hostility stems from his insight that contemporary astrologers were placing man ``back under the stars and the moon'', in effect saying that the ``external'' world of illusion was pre-determinant and ultimate. This represented a regress, or retardive, influence on the Tradition. Mouraviff speaks of the internal ``B-influences'' which properly signaled the birth of the inner consciousness (complete only if one could then discern both A \& C influence, and learn to use them consciously). The only way to graduate in this school or burden or task of God is to begin to transform the inner consciousness, both as regards itself, and as it regards the ``outer'' world. Hence the need for religion. On the contrary, both New Age \& Progress-through-technology represent precisely that ``externalization'' which religions of older times engaged in during a process or cycle of decay; \emph{man can only reach heaven by storm of the kingdom within}. Modernity hates religion so much because it sees in its caricature of such a direct competitor. This indicates occult influence. On the contrary, Ficino emphasizes a need for self-realization and transformation, rather than anything external.

\begin{quotex}
But we answer that on the contrary, the dignity of the heavens is taken away if one tries to interpret the movements of the heavens and the force of the signs according to the nature of the lowly animals. Rather, one should say that the heavens give something to them rather than receive from them.
\end{quotex}


The stars help activate seeds of the Logos, and relate to physical forms by being exemplars and catalysts. Being born under Taurus does not make one hardworking, patient, servile -- rather, as Ficino argues, it presents the soul born under such with certain possibilities, both good or bad. It is as the soul chooses in free will (a quality Ficino would agree \emph{is only present by orientation to the Good}) that it actualizes aspects, dimensions, and perfections of these possibilities. The insights are only potentially there; most lack the metaphysics to activate them. This is why Gornahoor's work is important.

\begin{quotex}
And rightly, because the greater material part conforms to the greater material cause, so- this is most important -- the formal part conforms to the formal causes, celestial and super-celestial.
\end{quotex}


Here we see Ficino acknowledging Being-Beyond-Being or the imperishable, which ultimately is potentially revealed in the rising spirit of the man who embraces personal alchemy. It is this man or woman who alone is ``free''; to embrace planetary predestination (or inner liberation) without \emph{soul-wisdom} of true free will is to fail to shape the potentials given to us ``under the moon'' \& to become pejoratively pagan. Ficino is showing us what Christianity's true aspiration should be -- neither to abolish paganism, nor to relapse into it.

\begin{quotex}
There are three kinds of foresight: firstly through divine infusion, secondly through natural instinct, thirdly through art. That which occurs through divine action includes divining through magical cult and the knowledge of daemons and spheres. That which is natural is assisted by a certain melancholic temperament or tempered at the same time by a heavenly power, by which the mind easily gathers into itself and divines through its own divinity, or by the observation of heavenly patterns. Art has four kinds: the most common is medicine, the most uncommon astrology. Physiognomy is closer to medicine, divinization to astrology.
\end{quotex}


This is not child's play nor amateur stuff.

\begin{quotex}
In all kinds of prediction judgement is very difficult: firstly because art is longer than a lifetime, but mainly because each art results, from the convergence of very many rules which are very different from each other, but each connected to the other. Secondly because the opportunities for judgement are limited, for neither is the mind always truthful nor do the things themselves exist in such a way that can be easily explained. Thirdly because experiential knowledge is faulty; since many factors can point to the same thing, it can scarcely be discerned by which one such or such an effect will proceed most powerful, and because experiments are not similar unless wholly similar in time and material. Fourthly because judgement itself is very difficult, not only for the same reasons, but also on account of the errors encroaching on the mind from the imagination and body. Here Isaiah says `tell me the things to come and I will call you gods'... lastly, through whatever art future things are questioned, they are foretold more completely from a certain gift of the soul, than through judgement. Here those who are unlearned in the art often judge more truthfully than those who are learned.
\end{quotex}


I want to note that the last emphasis (the same emphasis we found in the Pope's recent pronouncement on Pseudo-Dionysius) is not an invitation to cheap grace or modern ``Christendumb''. It is an emphasis on the virtue of Faith, which is predicated on childlike (not childish) Sincerity and Simplicity, just as the virtue of Hope is predicated on strength or fortitude, and the final virtue of Love is predicated on a successful passing of the ``steps'' or ``tests'' which are undergone in the formation of the second ``astral'' body. This is done under the union of the chakras through the heart, in subjection to the Lex Aeterna, which enables one to orient one's self in contact with the new reality uncovered in the second birth. This is the basis of alchemical Christianity, which was incarnated in Christian Europe, and which depended for its strength upon just such an awareness of Truth, even in the beginning stages. For better for worse, Christianity represented a collective cultural effort to make esotericism ``available'' for the masses without sacrificing the ideals. This the meaning of the Icon.

Hopefully, the above excerpts can serve to indicate a common direction in which Christians \& pagans alike can move which will be pleasing to the Creator, \& can alone lead to liberation into the freedom that all rational beings yearn towards, incidentally providing deliverance from the Brave New World by making us ``undigestible'' (since the goal of occult influence is precisely an inversion of this by ``processing'' everyone into the lowest common denominator, which can only be decay \& a shadow world of death). A good example of this for moderns might be Julian Lee\footnote{\url{http://www.celibacy.info/}}, who endeavors to unite Eastern meditation, Christian devotion, and awareness of astrological patterns in his geographical astrological work. However, those who rightly fear delusion will be wise to ``test all the spirits'' and to seek the wisdom that comes from God, since these are times in which even the elect can scarcely stand.

Excerpts taken from Ficino\footnote{\url{http://www.amazon.com/Marsilio-Ficino-Western-Esoteric-Masters/dp/1556435606}}, a book in the Esoteric series, edited by Angela Voss.



\flrightit{Posted on 2011-08-06 by Logres}

